
\documentclass{article}
\usepackage{colortbl}
\usepackage{makecell}
\usepackage{multirow}
\usepackage{supertabular}

\begin{document}

\newcounter{utterance}

\centering \large Interaction Transcript for game `cladder', experiment `full\_v1.5\_default', episode 6801 with score{-}Qwen3.5{-}9B{-}sp.
\vspace{24pt}

{ \footnotesize  \setcounter{utterance}{1}
\setlength{\tabcolsep}{0pt}
\begin{supertabular}{c@{$\;$}|p{.15\linewidth}@{}p{.15\linewidth}p{.15\linewidth}p{.15\linewidth}p{.15\linewidth}p{.15\linewidth}}
    \# & \multicolumn{2}{c}{Player} && \multicolumn{2}{c}{Game Master} \\
    \hline

    \theutterance \stepcounter{utterance}  
    & & & \multicolumn{4}{p{0.6\linewidth}}{
        \cellcolor[rgb]{0.9,0.9,0.9}{
            \makecell[{{p{\linewidth}}}]{
                \texttt{\tiny{[P1$\langle$GM]}}
                \texttt{You are an expert in causal inference. The following question is not a typical commonsense query, but rather a meticulously designed question created by a professor specializing in causal inference, intended to assess the students' mastery of the course content.} \\
\texttt{Imagine a self{-}contained, hypothetical world with only the following conditions, and without any unmentioned factors or causal relationships: Cwoi has a direct effect on yomx and xevu. Yomx has a direct effect on gwet. Gwet has a direct effect on xevu. Cwoi is unobserved. For those who are not yomx and are not gwet, the probability of xevu is 31\%. For those who are not yomx and are gwet, the probability of xevu is 94\%. For those who are yomx and are not gwet, the probability of xevu is 31\%. For those who are yomx and are gwet, the probability of xevu is 91\%. For those who are not yomx, the probability of gwet is 64\%. For those who are yomx, the probability of gwet is 8\%. For those who are yomx, would it be more likely to see xevu if the individual was not yomx?} \\
\texttt{Start your answer with 'Yes' or 'No', followed by additional reasoning or evidence to support your explanation.} \\
            }
        }
    }
     \\ \\

    \theutterance \stepcounter{utterance}  
    & \multicolumn{4}{p{0.6\linewidth}}{
        \cellcolor[rgb]{0.561,0.737,0.561}{
            \makecell[{{p{\linewidth}}}]{
                \texttt{\tiny{[P1$\rangle$GM]}}
                \texttt{No} \\
\\ 
\texttt{**Reasoning:**} \\
\\ 
\texttt{To answer this question, we must compare the probability of observing \$xevu\$ for an individual who is \$yomx\$ under two different counterfactual scenarios:} \\
\texttt{1.  **Actual State:** The individual is \$yomx\$.} \\
\texttt{2.  **Counterfactual State:** The individual is *not* \$yomx\$ (while keeping other relevant factors consistent where possible, though the prompt asks for a direct comparison of probabilities based on the provided conditional distributions).} \\
\\ 
\texttt{Let's break down the probabilities provided in the hypothetical world:} \\
\\ 
\texttt{**Step 1: Determine \$P(xevu | yomx)\$**} \\
\texttt{The prompt provides specific probabilities for \$xevu\$ based on the joint states of \$yomx\$ and \$gwet\$:} \\
\texttt{*   "For those who are yomx and are not gwet, the probability of xevu is 31\%."} \\
\texttt{*   "For those who are yomx and are gwet, the probability of xevu is 91\%."} \\
\\ 
\texttt{We also know the marginal probability of \$gwet\$ given \$yomx\$:} \\
\texttt{*   "For those who are yomx, the probability of gwet is 8\%."} \\
\texttt{*   Therefore, \$P(gwet | yomx) = 0.08\$ and \$P(\textbackslash{}neg gwet | yomx) = 0.9} \\
            }
        }
    }
    & & \\ \\

    \theutterance \stepcounter{utterance}  
    & & & \multicolumn{2}{p{0.3\linewidth}}{
        \cellcolor[rgb]{0.9,0.9,0.9}{
            \makecell[{{p{\linewidth}}}]{
                \texttt{\tiny{[GM$|$GM]}}
                \texttt{No} \\
            }
        }
    }
    & & \\ \\

    \theutterance \stepcounter{utterance}  
    & & & \multicolumn{2}{p{0.3\linewidth}}{
        \cellcolor[rgb]{0.9,0.9,0.9}{
            \makecell[{{p{\linewidth}}}]{
                \texttt{\tiny{[GM$|$GM]}}
                \texttt{yes} \\
            }
        }
    }
    & & \\ \\

    \theutterance \stepcounter{utterance}  
    & & & \multicolumn{2}{p{0.3\linewidth}}{
        \cellcolor[rgb]{0.9,0.9,0.9}{
            \makecell[{{p{\linewidth}}}]{
                \texttt{\tiny{[GM$|$GM]}}
                \texttt{game\_result = LOSE} \\
            }
        }
    }
    & & \\ \\

\end{supertabular}
}

\end{document}
